\section{Seguimiento}

Con el fin de evaluar el desempeño dinámico del controlador PII más allá del problema de regulación,
se realizaron experimentos de seguimiento en los que la referencia de posición varía en el tiempo.
Estas pruebas permiten analizar la interacción entre la acción de doble integración del controlador,
la no linealidad de zona muerta del actuador y el efecto atascamiento-deslizamiento modelado con
la fricción de Karnopp, todo ello en el intervalo de operación de $-3$ a $-2$ cm, en el que el actuador, limitado a $u\in[0;\,3{,}5]$~V, puede vencer la fricción estática en ambos sentidos.
El precedente experimental más cercano es el presentado por Alcántara et al.\ \cite{alcantara},
quienes demostraron en tiempo real tanto regulación como seguimiento para el mismo dispositivo
ECP-730 en configuración repulsiva (estable), empleando un esquema de linealización por
retroalimentación de estados combinado con un lazo externo con acción integral y control conmutado.
El presente trabajo extiende ese resultado al caso de configuración atractiva (inestable),
donde los puntos de equilibrio son inherentemente inestables y la tarea de seguimiento resulta
más exigente al requerir estabilización simultánea.

%=============================================================================
\subsection{Seguimiento de referencia senoidal}
%=============================================================================

La referencia de posición utilizada en esta prueba tiene la forma

\begin{equation}
  r(t) = x_0 + A\,\sin(2\pi f\, t),
  \label{eq:ref_senoidal}
\end{equation}

con valor nominal $x_0 = -2{,}5$ cm, amplitud $A = 0{,}5$ cm y frecuencia $f = 0{,}001$ Hz
(periodo $T = 1000$ s).
La amplitud fue seleccionada de modo que la referencia excursiona simétricamente entre
$-3$ cm y $-2$ cm, el mismo intervalo de la prueba de regulación. Se simularon dos periodos (2000~s).

\paragraph{Comportamiento esperado según el modelo lineal.}
Sin atascamiento-deslizamiento, el lazo cerrado lineal haría que la posición siguiera la referencia senoidal con un desfase pequeño y una atenuación de amplitud mínima. Si el sistema linealizado en lazo cerrado tiene función de transferencia $T(s)$, con $T(j\omega_0)=|T|\,e^{j\varphi}$, el error en estado estacionario ante una referencia senoidal de frecuencia $\omega_0 = 2\pi f$ es

\[
  e_{ss}(t) = A\bigl[1 - |T|\cos\varphi\bigr]\sin(\omega_0 t)
              - A\,|T|\sin\varphi\,\cos(\omega_0 t),
\]

expresión que muestra que el error es pequeño cuando $T(j\omega_0)$ es próximo a la unidad. Como se verá, con atascamiento-deslizamiento el error real lo domina el ciclo límite y no esta componente lineal.

\paragraph{Transiciones atascamiento-deslizamiento.}
Cada vez que la velocidad del levitante se anula, el sistema transita del régimen de deslizamiento al de atascamiento. Como se verá en los resultados, estas transiciones no ocurren solo cuando la referencia invierte su sentido ($t = T/4,\,3T/4,\ldots$), sino de forma continua, dando lugar a un ciclo límite.
El modelo de Karnopp detecta el atascamiento cuando se satisface la condición

\[
  |v| \leq DV = 0{,}02 \text{ cm/s},
\]

y la fuerza magnética neta no supera la fricción estática máxima $F_s = 20$~kg\,cm/s$^2$.
En esa condición, tanto la velocidad como la aceleración se anulan y el error de seguimiento
crece transitoriamente.
Una vez que la fuerza de control acumulada supera $F_s$ (\emph{breakaway}), el sistema
abandona el estado de reposo y retoma el seguimiento de la referencia.
La acción integral del controlador es la que, durante cada atascamiento, modifica el voltaje hasta que la fuerza neta supera la fricción estática; la comparación con un controlador PI se presenta al final de esta sección.

\begin{figure}[ht]
    \centering
    % Estilo común de las figuras de seguimiento (pgfplots).
\pgfplotsset{
  seguimiento/.style={
    width=0.92\textwidth, height=5.4cm,
    xlabel={tiempo [\si{s}]},
    grid=major, grid style={gray!25},
    tick label style={font=\small}, label style={font=\small}, legend style={font=\small},
    /pgf/number format/use comma,
    scaled ticks=false,
    y tick label style={/pgf/number format/fixed},
    x tick label style={/pgf/number format/fixed, /pgf/number format/1000 sep={}},
    table/col sep=comma,
    every axis plot/.append style={line join=round},
  },
}
\definecolor{tray}{RGB}{31,94,166}
\definecolor{refc}{RGB}{40,40,40}
\definecolor{ctl2}{RGB}{196,96,40}

\begin{tikzpicture}
\begin{axis}[seguimiento, name=principal, xmin=0, xmax=2000, ymin=-3.15, ymax=-1.85,
             xlabel={}, ylabel={posición [\si{cm}]},
             legend style={at={(0.99,0.03)}, anchor=south east}, legend cell align=left]
  \addplot[tray, line width=1pt] table[x=t, y=x] {images/seguimiento_results/tikz/sen_posicion.csv};
  \addlegendentry{posición $x_1$}
  \addplot[refc, densely dashed, line width=0.8pt] table[x=t, y=r] {images/seguimiento_results/tikz/sen_posicion.csv};
  \addlegendentry{referencia $r$}
  \fill[ctl2, opacity=0.15] (axis cs:244,-3.15) rectangle (axis cs:256,-1.85);
\end{axis}
\begin{axis}[seguimiento, at={(principal.south west)}, anchor=north west, yshift=-0.45cm,
             height=4.2cm, xmin=244, xmax=256,
             ylabel={detalle [\si{cm}]}, axis background/.style={fill=ctl2!4}]
  \addplot[tray, line width=1pt] table[x=t, y=x] {images/seguimiento_results/tikz/sen_detalle.csv};
  \addplot[refc, densely dashed, line width=0.8pt] table[x=t, y=r] {images/seguimiento_results/tikz/sen_detalle.csv};
\end{axis}
\end{tikzpicture}

    \caption[Seguimiento senoidal: posición]{Seguimiento de la referencia senoidal. Arriba, posición del levitante (línea continua) y referencia (discontinua) durante dos periodos; la banda sombreada marca la ventana de detalle de abajo (244--256~s, máximo de la referencia), donde se observa el ciclo límite de atascamiento-deslizamiento alrededor de la referencia.}
    \label{fig:sen_pos}
\end{figure}

La figura~\ref{fig:sen_pos} muestra que, a la escala de la referencia, la posición sigue a la senoide sin desfase ni atenuación apreciables, en concordancia con $|T(j\omega_0)|\approx 1$. La ventana de detalle revela, sin embargo, que el seguimiento no es continuo: el levitante avanza en saltos y permanece atascado entre ellos, formando un ciclo límite de aproximadamente $0{,}02$~cm pico a pico y periodo cercano a 1~s.

\FloatBarrier

\begin{figure}[ht]
    \centering
    % Estilo común de las figuras de seguimiento (pgfplots).
\pgfplotsset{
  seguimiento/.style={
    width=0.92\textwidth, height=5.4cm,
    xlabel={tiempo [\si{s}]},
    grid=major, grid style={gray!25},
    tick label style={font=\small}, label style={font=\small}, legend style={font=\small},
    /pgf/number format/use comma,
    scaled ticks=false,
    y tick label style={/pgf/number format/fixed},
    x tick label style={/pgf/number format/fixed, /pgf/number format/1000 sep={}},
    table/col sep=comma,
    every axis plot/.append style={line join=round},
  },
}
\definecolor{tray}{RGB}{31,94,166}
\definecolor{refc}{RGB}{40,40,40}
\definecolor{ctl2}{RGB}{196,96,40}

\begin{tikzpicture}
\begin{axis}[seguimiento, name=principal, xmin=0, xmax=2000, xlabel={}, ylabel={velocidad [\si{cm/s}]}]
  \addplot[tray, line width=0.6pt] table[x=t, y=v] {images/seguimiento_results/tikz/sen_velocidad.csv};
  \fill[ctl2, opacity=0.15] ({axis cs:244,0} |- {rel axis cs:0,0}) rectangle ({axis cs:256,0} |- {rel axis cs:0,1});
\end{axis}
\begin{axis}[seguimiento, at={(principal.south west)}, anchor=north west, yshift=-0.45cm,
             height=4.2cm, xmin=244, xmax=256, ylabel={detalle [\si{cm/s}]}, axis background/.style={fill=ctl2!4}]
  \addplot[tray, line width=0.8pt] table[x=t, y=v] {images/seguimiento_results/tikz/sen_detalle.csv};
\end{axis}
\end{tikzpicture}

    \caption[Seguimiento senoidal: velocidad]{Velocidad del levitante durante el seguimiento senoidal. Los pulsos corresponden a los eventos de deslizamiento y los tramos en cero al régimen de atascamiento (75\,\% del tiempo).}
    \label{fig:sen_vel}
\end{figure}

La figura~\ref{fig:sen_vel} muestra que la velocidad no evoluciona de forma suave: alterna pulsos de deslizamiento de hasta $\pm 1$~cm/s con tramos de velocidad nula. En los 2000~s simulados se registraron 4072 rupturas (unas dos por segundo) y el levitante permaneció atascado el 75\,\% del tiempo. El atascamiento, por tanto, no se limita a los instantes en que la referencia invierte su sentido, sino que ocurre de forma continua a lo largo de toda la prueba.

\FloatBarrier

\begin{figure}[ht]
    \centering
    % Estilo común de las figuras de seguimiento (pgfplots).
\pgfplotsset{
  seguimiento/.style={
    width=0.92\textwidth, height=5.4cm,
    xlabel={tiempo [\si{s}]},
    grid=major, grid style={gray!25},
    tick label style={font=\small}, label style={font=\small}, legend style={font=\small},
    /pgf/number format/use comma,
    scaled ticks=false,
    y tick label style={/pgf/number format/fixed},
    x tick label style={/pgf/number format/fixed, /pgf/number format/1000 sep={}},
    table/col sep=comma,
    every axis plot/.append style={line join=round},
  },
}
\definecolor{tray}{RGB}{31,94,166}
\definecolor{refc}{RGB}{40,40,40}
\definecolor{ctl2}{RGB}{196,96,40}

\begin{tikzpicture}
\begin{axis}[seguimiento, xmin=0, xmax=2000, ylabel={error $e=r-x_1$ [\si{cm}]}]
  \addplot[refc!60, thin] coordinates {(0,0) (2000,0)};
  \addplot[tray, line width=0.7pt] table[x=t, y=e] {images/seguimiento_results/tikz/sen_error.csv};
\end{axis}
\end{tikzpicture}

    \caption[Seguimiento senoidal: error]{Error de seguimiento $e(t)=r(t)-x_1(t)$ para la referencia senoidal.}
    \label{fig:sen_error}
\end{figure}

La figura~\ref{fig:sen_error} muestra que el error permanece acotado en una banda de $\pm 0{,}02$~cm (error máximo $0{,}0198$~cm, valor eficaz $0{,}0124$~cm) con media prácticamente nula. La amplitud de la banda varía ligeramente con la posición de la referencia, pero no se observan picos aislados: el error está dominado por el ciclo límite y no por la dinámica lineal de seguimiento.

\FloatBarrier

\begin{figure}[ht]
    \centering
    % Estilo común de las figuras de seguimiento (pgfplots).
\pgfplotsset{
  seguimiento/.style={
    width=0.92\textwidth, height=5.4cm,
    xlabel={tiempo [\si{s}]},
    grid=major, grid style={gray!25},
    tick label style={font=\small}, label style={font=\small}, legend style={font=\small},
    /pgf/number format/use comma,
    scaled ticks=false,
    y tick label style={/pgf/number format/fixed},
    x tick label style={/pgf/number format/fixed, /pgf/number format/1000 sep={}},
    table/col sep=comma,
    every axis plot/.append style={line join=round},
  },
}
\definecolor{tray}{RGB}{31,94,166}
\definecolor{refc}{RGB}{40,40,40}
\definecolor{ctl2}{RGB}{196,96,40}

\begin{tikzpicture}
\begin{axis}[seguimiento, name=principal, xmin=0, xmax=2000, ymin=-0.2, ymax=3.8, xlabel={},
             ylabel={señal de control [\si{V}]},
             legend style={at={(0.99,0.03)}, anchor=south east}, legend cell align=left, legend columns=3]
  \addplot[ctl2!45, line width=1.2pt] table[x=t, y=u_antes] {images/seguimiento_results/tikz/sen_control.csv};
  \addlegendentry{antes de la ZM}
  \addplot[tray, line width=0.5pt] table[x=t, y=u_desp] {images/seguimiento_results/tikz/sen_control.csv};
  \addlegendentry{después de la ZM}
  \addplot[red!70!black, densely dotted, line width=0.8pt] coordinates {(0,3.5) (2000,3.5)};
  \addlegendentry{$u_{\max}$}
  \fill[ctl2, opacity=0.15] (axis cs:244,-0.2) rectangle (axis cs:256,3.8);
\end{axis}
\begin{axis}[seguimiento, at={(principal.south west)}, anchor=north west, yshift=-0.45cm,
             height=4.2cm, xmin=244, xmax=256, ylabel={detalle [\si{V}]}, axis background/.style={fill=ctl2!4}]
  \addplot[tray, line width=0.8pt] table[x=t, y=u] {images/seguimiento_results/tikz/sen_detalle.csv};
\end{axis}
\end{tikzpicture}

    \caption[Seguimiento senoidal: señal de control]{Voltaje de control antes y después de la zona muerta durante el seguimiento senoidal. Abajo, detalle de 244 a 256~s: diente de sierra producido por la acción integral durante cada atascamiento.}
    \label{fig:sen_ctrl}
\end{figure}

La figura~\ref{fig:sen_ctrl} muestra que el voltaje sigue la forma de la referencia entre $1{,}33$ y $2{,}78$~V, siempre dentro del rango admisible $u \in [0;\,3{,}5]$~V y sin saturar. El detalle revela un perfil en diente de sierra: durante cada atascamiento la acción integral del PII hace crecer el voltaje hasta vencer la fricción estática, y tras la ruptura el voltaje cae bruscamente.

\FloatBarrier

%=============================================================================
\subsection{Seguimiento de referencia trapezoidal}
%=============================================================================

La referencia trapezoidal consiste en segmentos planos alternados en $x_1 = -2$ cm y
$x_1 = -3$ cm, conectados por rampas lineales de duración $250$ s cada una.
La señal describe así un perfil periódico de posición que combina fases de regulación
(segmentos planos) con fases de transición (rampas), constituyendo un escenario más exigente
que la referencia senoidal debido a la discontinuidad en la derivada de la referencia en los
instantes de inicio y fin de cada rampa.

\paragraph{Comportamiento durante la rampa.}
En la fase de rampa, la referencia varía a velocidad constante
$\dot{r} = (x_{f}-x_{i})/250 \text{ cm/s}$.
Según el modelo lineal, el error de seguimiento en esta fase converge a un valor constante determinado por el tipo del lazo.
En teoría, la segunda acción integral elimina el error de posición constante que un controlador de tipo uno mantendría durante la rampa; para la pendiente empleada, sin embargo, ese error sería del orden de $10^{-5}$~cm (véase la comparación al final de esta sección).

\paragraph{Comportamiento en los segmentos planos.}
Una vez que el levitante alcanza el nivel de referencia ($x_1 = -2$ cm o $x_1 = -3$ cm),
el sistema opera en modo de regulación.
Sin atascamiento-deslizamiento, la acción integral llevaría el error a cero; con él, como en la prueba de regulación, el error medio se anula pero el levitante oscila alrededor de la referencia en un ciclo límite.

\paragraph{Efecto de ruptura al inicio de cada rampa.}
El efecto de ruptura (\emph{breakaway}) se manifiesta al inicio de cada
rampa, cuando el sistema debe abandonar el estado de atascamiento en el que se encontraba
durante el segmento plano previo.
En ese instante, la fuerza magnética acumulada por la acción integral debe superar la
fricción estática máxima $F_s = 20$~kg\,cm/s$^2$ antes de que el levitante comience a desplazarse.
En la simulación, este retraso es de aproximadamente un segundo al inicio de la primera rampa y no produce un pico de error mayor que la banda del ciclo límite.

\paragraph{Señal de control durante la transición.}
En cada atascamiento, la acción integral del controlador modifica el voltaje hasta que la fuerza neta supera la fricción estática; tras la ruptura, el voltaje cae bruscamente, lo que produce un perfil en diente de sierra.
En todos los casos el voltaje se mantiene dentro del rango saturado $u \in [0;\,3{,}5]$ V,
sin alcanzar el límite superior.

\begin{figure}[ht]
    \centering
    % Estilo común de las figuras de seguimiento (pgfplots).
\pgfplotsset{
  seguimiento/.style={
    width=0.92\textwidth, height=5.4cm,
    xlabel={tiempo [\si{s}]},
    grid=major, grid style={gray!25},
    tick label style={font=\small}, label style={font=\small}, legend style={font=\small},
    /pgf/number format/use comma,
    scaled ticks=false,
    y tick label style={/pgf/number format/fixed},
    x tick label style={/pgf/number format/fixed, /pgf/number format/1000 sep={}},
    table/col sep=comma,
    every axis plot/.append style={line join=round},
  },
}
\definecolor{tray}{RGB}{31,94,166}
\definecolor{refc}{RGB}{40,40,40}
\definecolor{ctl2}{RGB}{196,96,40}

\begin{tikzpicture}
\begin{axis}[seguimiento, name=principal, xmin=0, xmax=2500, ymin=-3.15, ymax=-1.85,
             xlabel={}, ylabel={posición [\si{cm}]},
             legend style={at={(0.99,0.03)}, anchor=south east}, legend cell align=left]
  \addplot[tray, line width=1pt] table[x=t, y=x] {images/seguimiento_results/tikz/trap_posicion.csv};
  \addlegendentry{posición $x_1$}
  \addplot[refc, densely dashed, line width=0.8pt] table[x=t, y=r] {images/seguimiento_results/tikz/trap_posicion.csv};
  \addlegendentry{referencia $r$}
  \fill[ctl2, opacity=0.15] (axis cs:246,-3.15) rectangle (axis cs:266,-1.85);
\end{axis}
\begin{axis}[seguimiento, at={(principal.south west)}, anchor=north west, yshift=-0.45cm,
             height=4.2cm, xmin=246, xmax=266,
             ylabel={detalle [\si{cm}]}, axis background/.style={fill=ctl2!4}]
  \addplot[tray, line width=1pt] table[x=t, y=x] {images/seguimiento_results/tikz/trap_detalle.csv};
  \addplot[refc, densely dashed, line width=0.8pt] table[x=t, y=r] {images/seguimiento_results/tikz/trap_detalle.csv};
\end{axis}
\end{tikzpicture}

    \caption[Seguimiento trapezoidal: posición]{Seguimiento de la referencia trapezoidal. Arriba, posición y referencia durante dos periodos; abajo, detalle de 246 a 266~s, que incluye el inicio de la primera rampa en $t=250$~s.}
    \label{fig:trap_pos}
\end{figure}

La figura~\ref{fig:trap_pos} muestra que el levitante sigue las rampas y los segmentos planos sin error apreciable a la escala de la referencia. El detalle del inicio de la primera rampa muestra que el levitante permanece atascado alrededor de un segundo después de que la referencia comienza a moverse, y a partir de ahí sigue la rampa en escalones sucesivos de atascamiento-deslizamiento.

\FloatBarrier

\begin{figure}[ht]
    \centering
    % Estilo común de las figuras de seguimiento (pgfplots).
\pgfplotsset{
  seguimiento/.style={
    width=0.92\textwidth, height=5.4cm,
    xlabel={tiempo [\si{s}]},
    grid=major, grid style={gray!25},
    tick label style={font=\small}, label style={font=\small}, legend style={font=\small},
    /pgf/number format/use comma,
    scaled ticks=false,
    y tick label style={/pgf/number format/fixed},
    x tick label style={/pgf/number format/fixed, /pgf/number format/1000 sep={}},
    table/col sep=comma,
    every axis plot/.append style={line join=round},
  },
}
\definecolor{tray}{RGB}{31,94,166}
\definecolor{refc}{RGB}{40,40,40}
\definecolor{ctl2}{RGB}{196,96,40}

\begin{tikzpicture}
\begin{axis}[seguimiento, name=principal, xmin=0, xmax=2500, xlabel={}, ylabel={velocidad [\si{cm/s}]}]
  \addplot[tray, line width=0.6pt] table[x=t, y=v] {images/seguimiento_results/tikz/trap_velocidad.csv};
  \fill[ctl2, opacity=0.15] ({axis cs:246,0} |- {rel axis cs:0,0}) rectangle ({axis cs:266,0} |- {rel axis cs:0,1});
\end{axis}
\begin{axis}[seguimiento, at={(principal.south west)}, anchor=north west, yshift=-0.45cm,
             height=4.2cm, xmin=246, xmax=266, ylabel={detalle [\si{cm/s}]}, axis background/.style={fill=ctl2!4}]
  \addplot[tray, line width=0.8pt] table[x=t, y=v] {images/seguimiento_results/tikz/trap_detalle.csv};
\end{axis}
\end{tikzpicture}

    \caption[Seguimiento trapezoidal: velocidad]{Velocidad del levitante durante el seguimiento trapezoidal. El levitante permanece en reposo durante el primer segmento plano; a partir de la primera rampa, los pulsos de deslizamiento se mantienen durante el resto de la prueba.}
    \label{fig:trap_vel}
\end{figure}

La figura~\ref{fig:trap_vel} muestra que el levitante permanece en reposo durante el primer segmento plano, en el que parte del equilibrio. A partir de la primera rampa aparecen pulsos de deslizamiento de hasta $\pm 1$~cm/s, que continúan también en los segmentos planos posteriores: una vez iniciado, el ciclo límite no se extingue. En total se registraron 4400 rupturas en los 2500~s (dos periodos) y el levitante estuvo atascado el 78\,\% del tiempo.

\FloatBarrier

\begin{figure}[ht]
    \centering
    % Estilo común de las figuras de seguimiento (pgfplots).
\pgfplotsset{
  seguimiento/.style={
    width=0.92\textwidth, height=5.4cm,
    xlabel={tiempo [\si{s}]},
    grid=major, grid style={gray!25},
    tick label style={font=\small}, label style={font=\small}, legend style={font=\small},
    /pgf/number format/use comma,
    scaled ticks=false,
    y tick label style={/pgf/number format/fixed},
    x tick label style={/pgf/number format/fixed, /pgf/number format/1000 sep={}},
    table/col sep=comma,
    every axis plot/.append style={line join=round},
  },
}
\definecolor{tray}{RGB}{31,94,166}
\definecolor{refc}{RGB}{40,40,40}
\definecolor{ctl2}{RGB}{196,96,40}

\begin{tikzpicture}
\begin{axis}[seguimiento, xmin=0, xmax=2500, ylabel={error $e=r-x_1$ [\si{cm}]}]
  \addplot[refc!60, thin] coordinates {(0,0) (2500,0)};
  \addplot[tray, line width=0.7pt] table[x=t, y=e] {images/seguimiento_results/tikz/trap_error.csv};
\end{axis}
\end{tikzpicture}

    \caption[Seguimiento trapezoidal: error]{Error de seguimiento $e(t)=r(t)-x_1(t)$ para la referencia trapezoidal.}
    \label{fig:trap_error}
\end{figure}

La figura~\ref{fig:trap_error} muestra un error acotado en $\pm 0{,}022$~cm (valor eficaz $0{,}0114$~cm) y media nula, tanto en las rampas como en los segmentos planos. No se observan picos aislados al inicio de las rampas, sino una banda de amplitud casi uniforme debida al ciclo límite.

\FloatBarrier

\begin{figure}[ht]
    \centering
    % Estilo común de las figuras de seguimiento (pgfplots).
\pgfplotsset{
  seguimiento/.style={
    width=0.92\textwidth, height=5.4cm,
    xlabel={tiempo [\si{s}]},
    grid=major, grid style={gray!25},
    tick label style={font=\small}, label style={font=\small}, legend style={font=\small},
    /pgf/number format/use comma,
    scaled ticks=false,
    y tick label style={/pgf/number format/fixed},
    x tick label style={/pgf/number format/fixed, /pgf/number format/1000 sep={}},
    table/col sep=comma,
    every axis plot/.append style={line join=round},
  },
}
\definecolor{tray}{RGB}{31,94,166}
\definecolor{refc}{RGB}{40,40,40}
\definecolor{ctl2}{RGB}{196,96,40}

\begin{tikzpicture}
\begin{axis}[seguimiento, name=principal, xmin=0, xmax=2500, ymin=-0.2, ymax=3.8, xlabel={},
             ylabel={señal de control [\si{V}]},
             legend style={at={(0.99,0.03)}, anchor=south east}, legend cell align=left, legend columns=3]
  \addplot[ctl2!45, line width=1.2pt] table[x=t, y=u_antes] {images/seguimiento_results/tikz/trap_control.csv};
  \addlegendentry{antes de la ZM}
  \addplot[tray, line width=0.5pt] table[x=t, y=u_desp] {images/seguimiento_results/tikz/trap_control.csv};
  \addlegendentry{después de la ZM}
  \addplot[red!70!black, densely dotted, line width=0.8pt] coordinates {(0,3.5) (2500,3.5)};
  \addlegendentry{$u_{\max}$}
  \fill[ctl2, opacity=0.15] (axis cs:246,-0.2) rectangle (axis cs:266,3.8);
\end{axis}
\begin{axis}[seguimiento, at={(principal.south west)}, anchor=north west, yshift=-0.45cm,
             height=4.2cm, xmin=246, xmax=266, ylabel={detalle [\si{V}]}, axis background/.style={fill=ctl2!4}]
  \addplot[tray, line width=0.8pt] table[x=t, y=u] {images/seguimiento_results/tikz/trap_detalle.csv};
\end{axis}
\end{tikzpicture}

    \caption[Seguimiento trapezoidal: señal de control]{Voltaje de control antes y después de la zona muerta durante el seguimiento trapezoidal. Abajo, detalle de 246 a 266~s.}
    \label{fig:trap_ctrl}
\end{figure}

La figura~\ref{fig:trap_ctrl} muestra que el voltaje sigue el perfil trapezoidal entre $1{,}32$ y $2{,}78$~V, sin saturar. El detalle muestra cómo, al inicio de la rampa, el voltaje desciende de forma suave hasta que la fuerza neta vence la fricción estática, y a partir de ese instante adopta el mismo perfil en diente de sierra observado en la prueba senoidal.

\FloatBarrier

%=============================================================================
\subsection{Análisis comparativo: PI frente a PII}
%=============================================================================

Las pruebas de regulación y seguimiento muestran que el PII estabiliza el sistema, pero no permiten saber qué aporta la segunda acción integral respecto a un controlador con una sola.

Para contrastar el papel de la doble acción integral, ambas pruebas se repitieron con el controlador PI, con una sola acción integral, $C_{PI}(s)=178{,}14\,(s+15)(s+12{,}91)/[s\,(s+400)]$, en las mismas condiciones: escalón de $-2$ a $-3$~cm y un periodo de la referencia trapezoidal, con zona muerta, fricción de Karnopp y $u\in[0;\,3{,}5]$~V. La tabla~\ref{tab:pi_pii} resume los resultados.

\begin{table}[ht]
\centering
\caption{Comparación entre los controladores PI y PII con atascamiento-deslizamiento.}
\label{tab:pi_pii}
\begin{tabular}{|l|c|c|}
\hline
 & PI & PII \\ \hline
Sobrepaso ante el escalón & $25{,}9$\,\% & $23{,}6$\,\% \\ \hline
Ciclo límite pico a pico (escalón) & $0{,}031$~cm & $0{,}033$~cm \\ \hline
Duración media de cada atascamiento (escalón) & $0{,}26$~s & $0{,}36$~s \\ \hline
Duración media de cada atascamiento (trapezoidal) & $0{,}29$~s & $0{,}39$~s \\ \hline
Atascamiento más largo (trapezoidal) & $1{,}07$~s & $3{,}13$~s \\ \hline
Fracción del tiempo atascado (trapezoidal) & $74$\,\% & $76$\,\% \\ \hline
Error eficaz en las rampas & $0{,}0112$~cm & $0{,}0118$~cm \\ \hline
Error medio en las rampas & $-3\times10^{-5}$~cm & $1\times10^{-5}$~cm \\ \hline
\end{tabular}
\end{table}

Los resultados no muestran la ventaja que se esperaba de la segunda acción integral. El PII reduce ligeramente el sobrepaso, pero los atascamientos son en promedio más largos que con el PI y el ciclo límite tiene prácticamente la misma amplitud. En las rampas, el error residual que en teoría produce un controlador de tipo uno es $\dot r/K_v \approx 4\times10^{-5}$~cm para la pendiente empleada ($0{,}004$~cm/s), con $K_v=91{,}6$~s$^{-1}$ (sección~\ref{sec:analisis_ciclo}), muy inferior a la amplitud del ciclo límite, por lo que con esta referencia no es posible distinguir entre ambos controladores. Ambos operan sin saturar y con errores eficaces del mismo orden.

Para descartar que esta diferencia se deba a la sintonización particular de cada controlador, se evaluaron dos familias de controladores construidas a partir de los anteriores: en cada una se escaló la ganancia por un factor entre $0{,}5$ y $3$ y se desplazaron los ceros de baja frecuencia (los asociados a la acción integral) por un factor entre $0{,}25$ y $4$. De las combinaciones resultantes, 256 son estables con un margen de fase de al menos $30^\circ$ en todo el intervalo de $-3$ a $-2$~cm; de ellas, 54 PI y 51 PII completan, sin saturar ni perder la estabilidad, el escalón de $-2$ a $-3$~cm y una rampa de $-0{,}05$~cm/s, doce veces más pronunciada que la de la referencia trapezoidal. La figura~\ref{fig:r15_familias} muestra, para cada uno, la amplitud del ciclo límite y la duración media de los atascamientos.

\begin{figure}[ht]
    \centering
    % Generado a partir de calculus/Final_Bien/R15_sintonizacion (familias_R15.m, simular_R15.py).
\definecolor{tray}{RGB}{31,94,166}
\definecolor{ctl2}{RGB}{196,96,40}
\begin{tikzpicture}
\begin{axis}[width=0.85\textwidth, height=7cm, xmin=0, xmax=0.12, ymin=0, ymax=1.5,
    xlabel={ciclo límite pico a pico [\si{cm}]}, ylabel={duración media del atascamiento [\si{s}]},
    grid=major, grid style={gray!25}, tick label style={font=\small}, label style={font=\small},
    legend style={font=\small, at={(0.98,0.98)}, anchor=north east}, legend cell align=left,
    /pgf/number format/use comma, scaled ticks=false, x tick label style={/pgf/number format/fixed, /pgf/number format/precision=2},
    table/col sep=comma]
  \addplot[only marks, mark=o, mark size=2.2pt, ctl2] table[x=ciclo, y=atasc] {images/r15_results/tikz/familia_pi.csv};
  \addlegendentry{familia PI}
  \addplot[only marks, mark=square, mark size=2pt, tray] table[x=ciclo, y=atasc] {images/r15_results/tikz/familia_pii.csv};
  \addlegendentry{familia PII}
  \addplot[only marks, mark=*, mark size=3.5pt, ctl2, draw=black] coordinates {(0.0309,0.259)};
  \addlegendentry{PI de la comparación}
  \addplot[only marks, mark=square*, mark size=3.2pt, tray, draw=black] coordinates {(0.0339,0.352)};
  \addlegendentry{PII de la tesis}
\end{axis}
\end{tikzpicture}

    \caption[Familias PI y PII]{Ciclo límite y duración media de los atascamientos en el escalón de $-2$ a $-3$~cm para las familias de controladores PI y PII que operan sin saturar ($u\in[0;\,3{,}5]$~V). Los símbolos rellenos corresponden a los controladores de la tabla~\ref{tab:pi_pii}.}
    \label{fig:r15_familias}
\end{figure}

El ciclo límite más pequeño de la familia PII ($0{,}026$~cm) duplica el del mejor PI ($0{,}013$~cm), y para amplitudes de ciclo límite similares los PI se atascan durante menos tiempo. La amplitud del ciclo límite disminuye al aumentar la ganancia proporcional, no al añadir una segunda integración. La ventaja de la doble acción integral sí aparece en la rampa: con los controladores de la tabla~\ref{tab:pi_pii}, el error medio es de $0{,}13\times10^{-3}$~cm con el PII y de $0{,}56\times10^{-3}$~cm con el PI. Sin embargo, ambos son más de diez veces menores que el error eficaz que impone el ciclo límite en esa prueba, de unos $7{,}6\times10^{-3}$~cm y prácticamente igual con los dos controladores, por lo que no se traducen en un mejor seguimiento. La validación con integración \texttt{ode45} de los controladores extremos reproduce estos valores con diferencias menores al 8\,\%.

En síntesis, los experimentos de regulación y seguimiento demuestran que el controlador PII
diseñado es capaz de estabilizar el sistema de levitación magnética en configuración inestable
y de mantener un desempeño de seguimiento aceptable bajo los efectos combinados de la
zona muerta dinámica del actuador y la fricción de Karnopp, sin necesidad de compensación
explícita de ninguna de estas dos no linealidades.
